\section{Conclusion and future works}\label{sec:concl}

In this paper, we presented a new system for online, near real-time intrusion detection for monitoring Docker containers workloads. Our system can be implemented on the hosting machine where the containers are running, and it is based on the graph-based representation of the sequences of system calls generated by the containers. 
We tested the model against different types of anomalous behaviors, and the outcome of our experiments marks the ability to detect the majority of the attacks, despite no labelled attack sample was included in the training set. In addition, our results show that the three stage approach outperforms the two baselines built on one-class SVM and LOF models. Our system is flexible and may be trained on a greater number of expected hosted containers.

Future work will focus on the following: (i) evaluating the scalability of the intrusion detection system by considering an higher number of workloads in the training dataset;  (ii) extending the analysis to other container infrastructures, (iii) investigating the robustness and vulnerabilities of this system against adversarial attacks or other malicious behavior targeting the intrusion detection model; (iv) defining a distributed strategy for training the model and monitoring the containers that may migrate from one hosting machine to another.

\section{Acknowledgements}\label{sec:ack}
This research is partially funded by the EU H2020 ARCADIAN-IoT (Grant ID. 101020259) and partly by the H2020 CONCORDIA (Grant ID. 830927).